\documentclass[12pt, a4paper, oneside]{abntex2}

% Pacotes

\usepackage[utf8]{inputenc}
\usepackage[T1]{fontenc}
\usepackage[brazil]{babel}
\usepackage{times}
\usepackage{graphicx}
\usepackage{amsmath, amssymb, amsthm}
\usepackage{booktabs}
\usepackage{array}
\usepackage{float}
\usepackage{hyperref}
% \usepackage[alf, abnt-emphasize=bf]{abntex2cite}
\usepackage{microtype}
\usepackage{geometry}
\usepackage{indentfirst}
\usepackage{caption}
\usepackage{enumitem}
\usepackage{setspace}
\usepackage{footmisc}

\usepackage{amsmath}
\allowdisplaybreaks

\usepackage[style=abnt]{biblatex}
\addbibresource{referencias.bib}

\interfootnotelinepenalty=10000

% Configurações de página (ABNT NBR 14724)

\geometry{
  top=3cm,
  bottom=2cm,
  left=3cm,
  right=2cm
}

% ---
% Configurações de aparência dos títulos (Unificando a fonte)
% ---
\renewcommand{\ABNTEXchapterfont}{\rmfamily\bfseries}
\renewcommand{\ABNTEXsectionfont}{\rmfamily\bfseries}
\renewcommand{\ABNTEXsubsectionfont}{\rmfamily\bfseries}
\renewcommand{\ABNTEXsubsubsectionfont}{\rmfamily\bfseries}

% Opcional: Reduzir um pouco o tamanho exagerado do título do capítulo
\renewcommand{\ABNTEXchapterfontsize}{\Large}


% Espaçamento

\setlength{\parindent}{1.25cm}
\OnehalfSpacing


% Configurações do hyperref

\hypersetup{
  colorlinks=true,
  linkcolor=black,
  citecolor=black,
  urlcolor=black,
  pdftitle={Projeto de Pesquisa - Conectividade e Spillovers de Volatilidade no Sistema Bancário Brasileiro},
  pdfauthor={},
  pdfsubject={Projeto de Pesquisa - Monografia I},
  pdfkeywords={spillovers de volatilidade, risco sistêmico, sistema bancário brasileiro, Diebold-Yilmaz, teoria de redes}
  hidelinks
}

% Identificação

\titulo{Conectividade e Spillovers de Volatilidade no Sistema Bancário Brasileiro: Uma Análise de Redes via Diebold-Yilmaz}
\autor{Guilherme Henrique De Nardi}
\local{Osasco}
\data{2026}
\orientador{Diego Ferreira}
\instituicao{%
  Universidade Federal de São Paulo\\
  Escola Paulista de Política, Economia e Negócios
}
\tipotrabalho{Projeto de Pesquisa}

\preambulo{Projeto de pesquisa submetido como parte da avaliação da UC MONOGRAFIA I, do curso de graduação em Ciências Econômicas da Universidade Federal de São Paulo - Campus Osasco, sob orientação do Prof. Dr. Diego Ferreira.}

% ============================================================
\begin{document}

\frenchspacing

% ------------------------------------------------------------
% Elementos pré-textuais
% ------------------------------------------------------------
\pretextual

\imprimircapa

\imprimirfolhaderosto

% ---- Sumário ----
\pdfbookmark[0]{\contentsname}{toc}
\tableofcontents*
\cleardoublepage

% ============================================================
% Elementos textuais
% ============================================================
\textual

% ============================================================
\chapter{Introdução}
\label{cap:introducao}
% ============================================================

O sistema bancário brasileiro enfrentou, entre 2017 e 2026, um período marcado por choques
de naturezas distintas como uma pandemia global, o ciclo de aperto monetário do pós-pandemia,
e episódios de instabilidade idiossincrática como o caso Americanas. Esses episódios oferecem
um cenário propício para investigar como o risco se propaga entre instituições financeiras,
e se essa propagação muda de intensidade e direção conforme a natureza do choque. Risco sistêmico refere-se à possibilidade de que o colapso ou fragilização de uma instituição
se propague para o restante do sistema, ameaçando sua estabilidade como um todo. A literatura 
sobre redes financeiras estabelece que a arquitetura de interconexão que confere resiliência 
ao sistema em momentos de normalidade pode se converter em um mecanismo de amplificação dependendo
da magnitude de um choque \parencite{acemoglu2015}. Compreender
essa dinâmica, não apenas a magnitude mas também a evolução ao longo do tempo, é relevante 
tanto para a literatura quanto para a formulação de política macroprudencial. 

Mensurar a propagação de forma contínua e baseada em dados públicos é o objetivo central deste
trabalho, que investiga a rede de spillovers de volatilidade entre os principais bancos 
brasileiros ao longo do recorte 2017-2026. A pergunta que orienta a pesquisa é como a estrutura 
da rede entre essas instituições evolui ao longo do tempo e reage a eventos macroeconômicos
e idiossincráticos. A investigação busca não apenas identificar a existência de interdependência
entre as instituições, mas compreender mudanças na topologia da rede com variações no grau de 
conectividade e na importância relativa de cada instituição em diferentes contextos econômicos.

A hipótese do trabalho é de que a intensidade de spillovers se eleva em períodos identificáveis 
instabilidade financeira, como a pandemia do COVID-19, o ciclo de aperto monetário e episódios
idiossincráticos, refletindo maior interconectividade e potencial de propagações de choques, 
em linha com a dinâmica prevista pela literatura teórica de redes financeiras \parencite{acemoglu2015, allen2000}

A abordagem adotada baseia-se no modelo de Diebold e Yilmaz (2009, 2012, 2014), que utiliza um 
vetor autorregressivo (VAR) e a decomposição generalizada da variância dos erros de previsão 
(GFEVD) para mensurar a interdependência entre variáveis, construindo uma matriz de spillovers 
interpretável como a matriz de adjacência de uma rede complexa dirigida e ponderada. A variável-alvo é a 
volatilidade, não os retornos, estimada por um estimador de amplitude (range-based) que combina 
o componente de \textcite{rogers1991}, independente de tendência intradiária, com um termo 
de retorno overnight, capturando o fato de que decisões do Comitê de Política Monetária e 
resultados trimestrais são divulgados após o fechamento do pregão. A amostra é estratificada 
por tipo de controle de capital, porte e modelo de negócio, e organizada em dois painéis 
aninhados: um painel principal, com maior alcance temporal, e um painel de extensão que 
incorpora o Banco Pan, cuja relação de controle com o BTG Pactual permite testar diretamente 
se o vínculo societário entre instituições se traduz em maior comovimento de volatilidade.

Embora a metodologia de Diebold-Yilmaz tenha sido amplamente aplicada a mercados financeiros 
internacionais e a índices de ações de economias emergentes, sua aplicação ao nível de 
instituições bancárias individuais no Brasil permanece escassa. Os estudos brasileiros de risco 
sistêmico predominantemente dependem de dados de balanço ou de exposições interbancárias 
sigilosas, mantidas pelo Banco Central e inacessíveis à comunidade científica, o que compromete 
a comparabilidade dos resultados e restringe análises em escala ampla e frequência elevada. 
Este trabalho preenche essa lacuna em três frentes. Primeiro, estratifica a amostra por tipo de 
controle, modelo de negócio e porte, tornando a heterogeneidade institucional do sistema 
bancário brasileiro explícita na análise de spillovers. Segundo, compara explicitamente o ranking de 
importância sistêmica derivado dos spillovers com a segmentação prudencial do Banco Central do 
Brasil, testando se uma metodologia baseada inteiramente em dados públicos de mercado é capaz 
de aproximar a classificação regulatória das instituições consideradas sistemicamente relevantes.
Terceiro, incorpora métricas de teoria de redes além dos índices direcionais de primeira ordem 
propostos por Diebold e Yilmaz como centralidade de autovetor e PageRank, cuja formulação 
recursiva captura a posição estrutural de cada banco na topologia de contágio e testa a 
hipótese de vínculo societário por meio de uma regressão em painel de pares com efeitos fixos, 
abordagem mais robusta do que comparações agregadas por categoria nominal de controle.

O restante do trabalho está organizado em três capítulos. O Capítulo 2 revisa a literatura 
sobre risco sistêmico, conectividade em redes financeiras e a evolução metodológica do 
framework de Diebold-Yilmaz, situando as contribuições deste trabalho no debate. O Capítulo 3 
detalha a metodologia: a base de dados e os critérios de construção dos dois painéis amostrais, 
a estimação da volatilidade, a especificação do VAR e da GFEVD, o procedimento dinâmico via 
TVP-VAR, a regressão diádica e os procedimentos de visualização de redes. O Capítulo 4 
apresenta e discute os resultados, incluindo as considerações finais do trabalho.



% ============================================================
\chapter{Revisão Bibliográfica}
\label{cap:revisao}
% ============================================================

A noção de risco sistêmico na literatura econômica remonta a trabalhos
clássicos que destacam a fragilidade inerente das instituições financeiras.
O modelo de corridas bancárias de \textcite{diamond1983} estabelece que
bancos são estruturalmente vulneráveis a choques de confiança, nos quais a
retirada simultânea de depósitos pode levar ao colapso de instituições
solventes. Essa fragilidade não é somente individual, mas de natureza
sistêmica, na medida em que a falência de uma instituição pode desencadear
efeitos de contágio sobre o restante do sistema financeiro, configurando o
caráter sistêmico do risco bancário. A instabilidade pode
emergir de forma endógena relacionada à estrutura do próprio sistema, 
trabalhos subsequentes como o de \textcite{allen1998} reforçam essa perspectiva.
Embora esses modelos estabeleçam os fundamentos da fragilidade financeira,
ainda apresentam limitações por operarem sob hipóteses simplificadoras,
nas quais a heterogeneidade entre agentes e a complexidade das interações
financeiras são limitadas. Como resultado, esses modelos oferecem pouca
orientação sobre como a arquitetura das conexões financeiras influencia a
propagação de choques em sistemas contemporâneos.

A crise financeira global de 2008 representa um ponto de inflexão na pesquisa sobre
risco sistêmico em instituições financeiras. A importância da arquitetura de rede 
já havia sido demonstrada em \textcite{allen2000} de modo que as estruturas incompletas de 
reivindicações interbancárias se mostram mais propensas a episódios de contágio do que 
estruturas completas. Contudo, foi somente após o colapso de instituições de 
importância sistêmica que a fragilidade do sistema financeiro e a topologia da rede
financeira passaram a ocupar posição central no debate acadêmico e regulatório.

Por meio de uma análise bibliométrica extensa, \textcite{silva2017}
documentam essa virada: a literatura sobre estabilidade financeira evoluiu
para focar na topologia das redes interbancárias,
identificando o contágio via interconectividade como a principal fonte de
fragilidade sistêmica contemporânea. A interconectividade cria canais para o 
contágio e amplificação de choques no sistema financeiro, sendo este o ponto 
principal de \textcite{glasserman2014} que a identificam como um dos principais 
determinantes de crises financeiras recentes.

A mesma arquitetura de rede que promove resiliência em tempos de normalidade é a que propaga
contágio em momentos de estresse severo. Essa natureza dual é formalizada em 
\textcite{acemoglu2015}. O regime de propagação depende, portanto, da magnitude do choque 
inicial. 
Uma rede altamente conectada atua como um mecanismo de absorção de choques para
choques de baixa magnitude, mas se converte em mecanismo de amplificação e
propagação sistêmica quando os choques ultrapassam um limiar crítico. Essa dinâmica de limiar
é formalizada de maneira mais ampla pela literatura de cascatas em redes financeiras.
Os mecanismos pelos quais as perdas se propagam em cascata é mapeado extensamente em 
\textcite{glasserman2016},
distinguindo canais de contágio diretos e indiretos. A 
probabilidade e magnitude de uma cascata dependem conjuntamente da topologia de 
rede e da distribuição de exposições, fortalecendo a percepção de que a mensuração 
empírica de conectividade é condição necessária para a avaliação do risco sistêmico.

A literatura converge na ideia de que a interconectividade é elemento central
para a compreensão do risco sistêmico. Diverge, porém, quanto à sua mensuração
empírica, especialmente em contextos nos quais a estrutura da rede não é diretamente
observável. Como resultado, grande parte da literatura depende de proxies indiretas que
aproximam imperfeitamente a topologia da rede, ou de bases granulares de acesso restrito
mantidas por reguladores, o que compromete a comparabilidade dos resultados e restringe
análises empíricas em escala ampla e frequência elevada.

O problema da mensuração indireta encontra uma solução metodológica no spillover index, 
proposto por \textcite{diebold2009} a partir de dados de mercado de acesso público. 
A metodologia utiliza um modelo de vetores autorregressivos (VAR) e a decomposição da 
variância dos erros de previsão para gerar um índice agregado de conectividade, mensurando a 
interdependência entre ativos ou mercados financeiros. 
Em outras palavras, o índice quantifica a parcela da variância de
uma variável atribuível a choques nas demais, em contraposição à parcela atribuível à
própria. Por meio da representação MA($\infty$) do VAR e da decomposição de Cholesky,
os autores decompõem a variância do erro de previsão em parcelas de variância própria
(\textit{own variance shares}), decorrentes de choques na própria variável, e parcelas
de variância cruzada (\textit{cross variance shares}). O \textit{spillover index} é
construído pela soma das contribuições de choques originados em $i$ sobre o erro de
previsão de $j$ (para todo $i \neq j$), normalizada pela variância total do sistema.

Entretanto, o \textit{spillover index} original apresentava duas limitações
principais sendo a sensibilidade à ordem das variáveis na decomposição de variância e a
dificuldade em identificar a origem e o destino dos choques. Para superar esses
desafios, uma refinação da metodologia foi apresentada em \textcite{diebold2012} ao adotar 
a Decomposição Generalizada da Variância dos Erros de Previsão (GFEVD), proposta 
originalmente por \textcite{koop1996} e \textcite{pesaran1998}. Ao adotar a GFEVD, os autores
eliminaram o problema inerente à ordenação das variáveis, característica da decomposição
de Cholesky, possibilitando uma identificação invariante da topologia de transmissão de
risco. O refinamento metodológico permitiu, ainda, a introdução do conceito de
\textit{net spillover}, definido como a diferença entre a contribuição que uma variável
transmite ao sistema e aquela que dela recebe, classificando cada variável como
transmissora ou receptora líquida de volatilidade. 
O balanço histórico de \textcite{diebold2023} consolida o modelo de conectividade dinâmica
como a base para a análise macroprudencial, demonstrando que os \textit{spillovers} de mercado
refletem fielmente a topologia de transmissão do risco sistêmico.

A matriz de conectividade resultante da GFEVD admite, ainda, uma leitura sob a ótica 
da teoria das redes como a matriz de adjacência de uma rede complexa,
intuição formalizada em \textcite{diebold2014}. Diferentemente da maior parte da literatura 
aplicada de redes financeiras, que analisa apenas a existência de conexões entre instituições, 
a estrutura derivada da GFEVD configura um grafo direcionado e ponderado. 
A natureza direcionada da rede captura a dinâmica de assimetria do risco financeiro, dado 
que a contribuição de $i$ para a variância de $j$ não coincide, em geral, com a 
contribuição inversa.

Dessa estrutura, \textcite{diebold2014} derivam métricas de conectividade com
correspondência direta nas medidas fundamentais da teoria de grafos.
O conceito de \textit{from-degree}, que mede a conectividade total recebida pelo nó
(equivalente ao \textit{in-degree} em teoria dos grafos), funciona como indicador de
vulnerabilidade da variável a choques transmitidos pelos demais elementos do sistema.
Em contrapartida, o \textit{to-degree}, equivalente ao \textit{out-degree}, mede a
conectividade transmitida pelo nó aos demais e identifica os elementos com maior potencial
de propagação sistêmica. O índice agregado de conectividade do sistema, por sua vez,
sintetiza a densidade da interdependência da rede, sendo equivalente conceitualmente ao
grau médio normalizado de um grafo ponderado.

A abordagem de janelas móveis (\textit{rolling-window}), embora permita capturar a
evolução temporal da conectividade, impõe uma limitação relevante em que os parâmetros do
modelo são recalculados em janelas fixas e discretas, o que pode obscurecer transições
graduais na estrutura da rede. A alternativa é estimar os parâmetros de forma contínua,
sem a definição de uma janela fixa. 
É o que fazem \textcite{antonakakis2020} e \textcite{chatziantoniou2021} 
ao propor a estimação do índice de conectividade via
modelos de parâmetros variantes no tempo (\textit{Time-Varying Parameter VAR} - TVP-VAR).
Essa extensão metodológica preserva a interpretabilidade do
\textit{framework} de Diebold-Yilmaz, ao mesmo tempo em que oferece uma representação
mais contínua e sensível da dinâmica de transmissão de risco, tornando-se especialmente
adequada para a análise de sistemas financeiros sujeitos a choques abruptos e regimes
variantes. A escolha entre rolling-window e TVP-VAR envolve um \textit{trade-off} entre
simplicidade interpretativa e suavidade na captura de transições de regime. A abordagem 
de janelas móveis, embora consiga atenuar efeitos de memória longa característicos de séries de 
volatilidade financeira não é suficiente para identificar episódios de elevação da conectividade
associados já que data defasadamente os movimentos de acordo com o tamanho especificado da janela.
Por esse motivo, o TVP-VAR é adotado como especificação principal neste trabalho, com a janela 
móvel mantida apenas como comparação com a convenção da literatura de referência.

A representação espectral da decomposição da variância permite também analisar os spillovers 
no domínio da frequência. O risco sistêmico tem dinâmicas distintas. A conectividade gerada em 
altas frequências reflete o processamento rápido de informações no curto prazo em momentos de 
calmaria, enquanto a conectividade gerada em baixa frequências sinaliza choques persistentes e 
mudanças estruturais nas expectativas dos agentes, sendo esse o ponto central de 
\textcite{barunik2018}. O modelo exige um horizonte de aproximação e uma extensão amostral 
longos para que o algoritmo de Fourier mapeie as baixas frequências com precisão. 
Essas escolhas introduzem o risco de sofrer com \textit{overfitting} e problemas de 
graus de liberdade, especialmente em janelas amostrais mais curtas. 

Também é possível estender a abordagem para o domínio
quantílico conforme proposto em \textcite{ando2022}. Os autores demonstraram empiricamente que 
a transmissão de choques e o risco de contágio são fenômenos de cauda, se apresentam virtualmente
ausentes na mediana e se expandem de forma severa em períodos de estresse. Dada a natureza assimétrica
e não-linear da conectividade financeira, os modelos baseados na média condicional tendem a subestimar
o potencial de contágio sistêmico em momentos de crise. A estimação depende de dados de densidade
alta e amostras maiores, o que inviabiliza sua aplicação. A dinâmica temporal de janelas móveis
consegue capturar a transição do sistema para a cauda uma vez que a média condicional 
é dominada por observações extremas em momento de crise, além de ser computacionalmente acessível
e estável.  

O \textit{spillover index} foi utilizado para analisar a volatilidade dos retornos da rede de 
instituições financeiras norte-americanas e europeias no período 2004--2014 e caracterizar a 
dinâmica da conectividade ao longo da crise financeira de 2008 em \textcite{diebold2016}. 
Os autores argumentam que, no período pós-2008, a transmissão de
volatilidade se tornou mais bidirecional, com aumento expressivo da influência das
instituições europeias sobre o mercado norte-americano. Além disso, os autores demonstram
que, após o colapso do Lehman Brothers, a crise financeira teria adquirido caráter
efetivamente global. O estudo evidencia o aprofundamento das conexões entre instituições
financeiras em períodos de estresse, corroborando a interpretação de
\textcite{acemoglu2015} de que a interconectividade que estabiliza em normalidade 
amplifica choques sob estresse severo.

Para os mercados das Américas \textcite{diebold2011} analisam empiricamente os 
\textit{spillovers} de volatilidade entre cinco mercados de ações: EUA, México, Chile, 
Argentina e Brasil. Os autores identificam que, embora as
economias centrais dominem a transmissão em momentos de estresse agudo, as economias
latino-americanas exibem fluxos de transmissão regional não desprezíveis. Para o Brasil,
especificamente, os resultados indicam que o país não opera como receptor passivo de
volatilidade, mas atua como \textit{hub} regional de propagação de volatilidade,
refletindo sua centralidade econômica na América Latina. Esse achado motiva investigações
mais aprofundadas sobre como a posição do Brasil na rede de \textit{spillovers} se
comporta em períodos posteriores, marcados por choques idiossincráticos domésticos.

O risco soberano oferece outro campo de aplicação do \textit{framework} de Diebold-Yilmaz, 
mensurando os \textit{spillovers} de volatilidade de \textit{credit default swaps} 
(CDS) entre economias emergentes e avançadas, foco central em \textcite{bostanci2020}.
O estudo evidencia que a rede de spillovers é altamente integrada e que os mercados emergentes
(incluindo o Brasil) assumiram um papel central como transmissores líquidos de volatilidade
financeira no período pós-crise. A unidade de análise, porém, é o risco soberano agregado.
O CDS reflete o risco de crédito do país como um todo e não a estrutura de transmissão 
entre instituições individualmente. Permanece em aberto se essa centralidade do Brasil 
como transmissor líquido se reproduz ou se inverte quando a rede é decomposta 
nas instituições que a compõem. O país pode ser transmissor líquido agregado 
enquanto, internamente, alguns bancos atuam como receptores.


O mercado brasileiro enfrentou, no entanto, ao longo de 2019-2025, 
uma combinação de choques de natureza distinta daquela observada nos sistemas avançados 
analisados em \textcite{diebold2016}. Uma pandemia global, um ciclo abrupto de aperto monetário
e episódios de instabilidade como a fraude contábil da Americanas sugerindo que a estrutura de rede de 
\textit{spillovers} brasileira possui determinantes distintos daqueles observados nas
economias avançadas. O mercado interbancário brasileiro
apresenta estrutura concentrada. As primeiras evidências empíricas em \textcite{cajueiro2008} 
de topologia do mercado interbancário brasileiro sob a ótica das redes complexas
demonstram que o sistema exibe elevada heterogeneidade estruturada em torno de
\textit{money centers}, em que o porte e o tipo de controle do capital determinam a
centralidade da instituição na rede.

Evidências para o caso brasileiro também sugerem que a estrutura do sistema interbancário brasileiro
apresenta uma elevada concentração de conectividade em um número reduzido de instituições 
de importância sistêmica. Em \textcite{silva2016}, os autores
caracterizam a rede interbancária brasileira por meio de medidas de teoria de redes. O
artigo documenta que o mercado interbancário brasileiro possui uma estrutura de
\textit{core-periphery}, em que o núcleo é composto por um pequeno grupo de bancos
grandes e altamente conectados, e a periferia, por um grande número de bancos menores
que se conectam ao núcleo, mas raramente entre si próprios. O estudo, contudo, depende
de dados granulares e sigilosos mantidos pelo Banco Central do Brasil, inacessíveis ao
público e à comunidade científica.

O mesmo problema de dependência de dados sigilosos aparece em \textcite{cont2013}, que
propõem o \textit{Contagion Index}, a perda esperada do sistema condicional à inadimplência 
de uma instituição em um período de estresse macroeconômico a partir de uma base de exposições
cedidas pelo Banco Central do Brasil. Os achados dialogam com \textcite{silva2016} já que
o sistema exibe topologia \textit{scale-free}, em que a maioria dos nós possui poucas conexões, 
enquanto alguns nós, conhecidos como \textit{hubs}, concentram um grande número de conexões. 
Sob essa topologia, a concentração da exposição e do risco sistêmico em poucos nós centrais
favorece a propagação de perdas em cascata, com choques macroeconômicos atuando como
mecanismo central de amplificação. A pesquisa, no entanto, tem seu mecanismo de transmissão
modelado exclusivamente pelo balanço, ignorando o canal dos preços de mercado pelo qual a
incerteza sobre uma instituição se transmite às demais antes que qualquer exposição materialize
a perda. 

O $\Delta$CoVaR aplicado ao caso brasileiro por \textcite{moura2024}, com dados de mercado,
demonstra que a contribuição marginal para o risco sistêmico no Brasil está ligada ao perfil de
tomada de risco das instituições. Ainda que o estudo evidencie empiricamente
que a vulnerabilidade de balanços é relevante para o risco sistêmico, o método não captura
a direcionalidade e a topologia da rede em tempo real. Embora o $\Delta$CoVaR
identifique a magnitude do impacto potencial de uma instituição sobre o sistema, ele falha
em mapear os canais de transmissão e o efeito cascata das interações, tornando necessária 
uma abordagem complementar baseada em spillovers que desenhe a arquitetura das interações.

Em síntese, as evidências disponíveis para o caso brasileiro partem, predominantemente,
de dados sigilosos do Banco Central do Brasil. Essa dependência cria uma lacuna metodológica
que exige abordagens complementares fundamentadas em dados públicos de mercado, como os
\textit{spillovers} de volatilidade capturados pela metodologia de Diebold-Yilmaz, que
permitam uma análise dinâmica, transparente e em alta frequência da conectividade
sistêmica, dispensando o acesso a bases sigilosas.

A lacuna de dados públicos encontra resposta parcial em \textcite{silva2025}, que visa 
resolver investigando a conectividade através de dados de ações bancárias listadas no Ibovespa. 
Utilizando um modelo TVP-VAR e complementando o estudo com métricas de redes complexas,
os autores buscam entender como o contágio financeiro se propaga através da rede brasileira. 
O estudo demonstra que a estrutura
de transmissão de volatilidade entre bancos brasileiros é instável e sensível a choques 
macroeconômicos e políticos, atingindo picos superiores a $50\%$ em períodos de instabilidade
política e econômica com Bradesco, Banco do Brasil e Itaú operando como os nós centrais
na transmissão de choques. Contudo, o estudo define seu universo amostral por critério de
liquidez, sem estratificar as instituições, o que limita a capacidade de caracterizar como a 
heterogeneidade institucional do sistema bancário influencia a estrutura de transmissão
de risco.

O presente trabalho propõe-se a preencher essas lacunas em três frentes. 
Primeiro, estratifica a amostra por tipo de controle (público, privado nacional e estrangeiro), modelo de negócio (varejo, investimento, atacado e consumo) e porte, tornando a 
heterogeneidade institucional do sistema bancário brasileiro explícita na análise de 
\textit{spillovers}. Segundo, vai além das métricas reportadas por \textcite{silva2025} 
ao comparar explicitamente o ranking de importância sistêmica derivado dos \textit{spillovers}
com a segmentação prudencial de D-SIBs publicada pelo Banco Central do Brasil, testando se uma 
metodologia baseada em dados públicos de mercado é capaz de aproximar a classificação 
regulatória. Terceiro, incorpora métricas de redes ampliadas, \textit{eigenvector centrality}, 
PageRank e regressão diádica, sendo esta última raramente reportada em estudos de 
\textit{spillovers} brasileiros e particularmente adequada para testar se há 
\textit{clustering} na rede de transmissão de risco.

\chapter{Metodologia}
\label{cap:metodologia}

Esta seção apresenta as escolhas metodológicas do trabalho em três subseções. A
Subseção~\ref{subsec:dados} descreve a base de dados, os critérios de seleção da amostra,
a construção das séries de volatilidade e os procedimentos de tratamento dos dados. A
Subseção~\ref{subsec:dy} apresenta a topologia de redes via Diebold-Yilmaz, incluindo a
especificação do VAR, a identificação por decomposição generalizada da variância e os
índices de spillover. A Subseção~\ref{subsec:redes} detalha a visualização de redes e a
interpretação econômica das métricas de centralidade. A notação é declarada na
Subseção~\ref{subsec:dados} e mantida ao longo de toda a seção. 
%TODO
%Os pacotes computacionais
%utilizados, as versões e o código completo encontram-se no Apêndice Computacional.

% ──────────────────────────────────────────────────────────────────────────────
\section{Base de Dados}
\label{subsec:dados}
% ──────────────────────────────────────────────────────────────────────────────

\subsubsection*{Seleção da amostra}

A amostra é composta por oito instituições bancárias com ações negociadas na B3 entre
fevereiro de 2019 e dezembro de 2025. A composição foi definida de modo a maximizar três
dimensões de heterogeneidade relevantes para o estudo de risco sistêmico: (i)~\textit{tipo
de controle}, abrangendo bancos privados nacionais, público federal, público estadual e
privados estrangeiros; (ii)~\textit{perfil de negócio}, cobrindo varejo massificado, banco
de investimento, atacado corporativo e crédito ao consumo; e (iii)~\textit{porte},
incluindo tanto \textit{large caps} quanto \textit{mid caps}.

A amostra contempla as oito instituições bancárias com ações listadas na B3 que reúnem, em
conjunto, representatividade sistêmica e diversidade de perfis de controle e modelo de negócio.
Itaú Unibanco (ITUB4) e Bradesco (BBDC4) figuram como os dois maiores bancos privados nacionais,
funcionando como \textit{hubs} centrais da rede de conectividade. O Banco do Brasil (BBAS3),
como maior instituição pública federal, captura a exposição do sistema a políticas governamentais.
O Santander (SANB11) representa o principal canal de transmissão de choques internacionais, 
na condição de maior banco estrangeiro do país. O BTG Pactual (BPAC11), maior banco de 
investimento independente da América Latina, contribui com um perfil de \textit{wholesale} 
distinto dos bancos de varejo tradicionais. O Banrisul (BRSR6), único banco público estadual 
com liquidez suficiente para o estudo, introduz uma dimensão regional (Rio Grande do Sul) à 
amostra. O ABC Brasil (ABCB4), por sua vez, diversifica tanto o tipo de controle 
(privado estrangeiro) quanto o modelo de negócio (\textit{wholesale} puro). Por fim, o 
Banco Pan (BPAN4) representa o segmento de bancos médios de varejo, com a particularidade de 
manter relação societária com o BTG Pactual, o que pode gerar \textit{spillovers} 
par-a-par relevantes entre as duas instituições.

Três instituições com presença relevante no mercado de capitais foram deliberadamente
excluídas. O Banco BMG (BMGB4) apresenta perfil de negócio redundante com o Banco Pan,
ambos concentrados em crédito consignado e varejo, além de liquidez em bolsa
materialmente inferior. A B3~S.A.\ (B3SA3) não é uma instituição bancária, mas uma
infraestrutura de mercado; sua inclusão alteraria a natureza do estudo. BB Seguridade
(BBSE3) e Porto Seguro (PSSA3) são seguradoras e, portanto, sua inclusão deslocaria o foco
do sistema bancário para uma rede financeira mais ampla, comprometendo a comparabilidade
com as Instituições Financeiras Sistemicamente Importantes Domésticas (D-SIBs) designadas
pelo Banco Central do Brasil (BCB).

O recorte temporal de cada instituição não é uma escolha discricionária, mas
decorre diretamente do critério de liquidez descrito na Subseção~\ref{sec:tratamento-dados}
(volume financeiro diário inferior a R\$~100 mil). Seis das oito instituições —
ITUB4, BBDC4, BBAS3, SANB11, BRSR6 e ABCB4 — nunca ficam abaixo desse piso ao
longo de toda a base disponível, desde 2014. As duas exceções definem as
fronteiras efetivas da amostra. O BPAC11 só passa a negociar de forma
continuamente líquida a partir de 18/08/2017 — antes disso, a ação, listada em
fevereiro de 2017, ainda está em fase de formação de \textit{book}. O BPAN4, por sua vez,
só se torna continuamente líquido a partir de 16/01/2019: no terceiro trimestre
de 2018, por exemplo, o papel permanece abaixo do piso de liquidez em 44 dos 63
pregões do período, com média de apenas 48 negócios diários. 

Essas duas fronteiras — determinadas pelos dados, não escolhidas a priori — geram
naturalmente dois painéis aninhados, apresentados na Tabela~\ref{tab:paineis}.
O painel principal (N~=~7) exclui o Banco Pan e cobre o período de outubro de
2017 a maio de 2026, alcançando 2.148 pregões ao longo de 8,7 anos — o maior
alcance temporal possível dado o critério de liquidez das sete instituições
restantes. O painel de extensão (N~=~8) inclui o Banco Pan e cobre janeiro de
2019 a outubro de 2025, totalizando 1.690 pregões ao longo de 6,8 anos.

\begin{table}[htb]
\centering
\caption{Painéis amostrais com fronteiras determinadas pelos dados}
\label{tab:paineis}
\begin{tabular}{lccc}
\toprule
Painel & Janela & Observações & Papel \\
\midrule
N = 7 (sem Banco Pan) & 10/2017 a 05/2026 & 2.148 & Principal \\
N = 8 (com Banco Pan) & 01/2019 a 10/2025 & 1.690 & Extensão \\
\bottomrule
\end{tabular}
\legend{Fonte: elaboração própria, a partir de ComDinheiro.}
\end{table}


\subsection*{6.1.2 Construção das séries de volatilidade}

\subsubsection*{Justificativa da variável-alvo}

A variável-alvo deste estudo é a volatilidade das ações dos bancos selecionados, e não os
retornos. Essa escolha é motivada pelo próprio objeto de análise: importância sistêmica, no
sentido adotado pelo BCB para a designação de D-SIBs, é um conceito de segundo momento e
refere-se à capacidade de uma instituição propagar \textit{estresse} pelo sistema
financeiro, não apenas retornos. Spillovers de volatilidade capturam precisamente essa
dimensão e são tipicamente maiores, mais informativos e mais robustos a mudanças de regime
do que spillovers de retornos \parencite{diebold2009}. Ademais, a janela amostral escolhida
(2019--2025) concentra episódios cujos efeitos se manifestaram primordialmente na
volatilidade como COVID-19, ciclo de aperto da Selic
em 2022 e fraude contábil da Americanas em 2023. Por fim, a presença de \textit{mid caps}
na amostra, cujo volume de negociação é irregular, recomenda o uso de medidas
\textit{range-based}, mais robustas ao ruído de microestrutura do que retornos diários
\textit{close-to-close}.

\subsubsection*{Estimador de Yang-Zhang}

A volatilidade é mensurada pelo estimador \textit{range-based} de
\textcite{yang2000}, que combina três componentes em uma média ponderada que minimiza a
variância do estimador:

\begin{equation}
  \hat{\sigma}^{2}_{\mathrm{YZ}}
    = \hat{\sigma}^{2}_{o}
    + k\,\hat{\sigma}^{2}_{c}
    + (1-k)\,\hat{\sigma}^{2}_{\mathrm{RS}},
  \label{eq:yz}
\end{equation}

\noindent onde, para uma janela de $n$ dias úteis:

\begin{align}
  \hat{\sigma}^{2}_{o}
    &= \frac{1}{n-1}\sum_{t=1}^{n}
       \bigl[\log(O_t/C_{t-1}) - \bar{r}_o\bigr]^{2},
  \label{eq:overnight}\\[4pt]
  \hat{\sigma}^{2}_{c}
    &= \frac{1}{n-1}\sum_{t=1}^{n}
       \bigl[\log(C_t/O_t) - \bar{r}_c\bigr]^{2},
  \label{eq:openclose}\\[4pt]
  \hat{\sigma}^{2}_{\mathrm{RS}}
    &= \frac{1}{n-1}\sum_{t=1}^{n}
       \Bigl[
         \log(H_t/C_t)\log(H_t/O_t)
         + \log(L_t/C_t)\log(L_t/O_t)
       \Bigr],
  \label{eq:rs}\\[4pt]
  k &= \frac{0{,}34}{1{,}34 + \dfrac{n+1}{n-1}}.
  \label{eq:k}
\end{align}

\noindent Aqui $H_t$, $L_t$, $O_t$ e $C_t$ denotam os preços máximo, mínimo, de abertura
e de fechamento no dia $t$, respectivamente. O parâmetro $k$ em~\eqref{eq:k} decorre da
minimização analítica da variância do estimador conjunto \parencite{yang2000}. Os
termos $\bar{r}_o = n^{-1}\sum_{t}\log(O_t/C_{t-1})$ e
$\bar{r}_c = n^{-1}\sum_{t}\log(C_t/O_t)$ são as médias amostrais dos retornos
\textit{overnight} e \textit{open-to-close} dentro da janela, respectivamente.

A janela é fixada em $n = 21$ dias úteis, alinhando a medida ao conceito de volatilidade
mensal. A variável efetivamente utilizada nas estimações é
$y_t = \log\hat{\sigma}_{\mathrm{YZ},t}$, onde
$\hat{\sigma}_{\mathrm{YZ},t} = \sqrt{\hat{\sigma}^{2}_{\mathrm{YZ},t}}$. Essa
transformação logarítmica é padrão na literatura de spillovers \parencite{diebold2012} e
aproxima a distribuição dos choques de um padrão gaussiano, facilitando a inferência no
VAR.

\subsubsection*{Por que Yang-Zhang e não estimadores alternativos?}

O estimador de Garman-Klass ignora a variação \textit{overnight} ($C_{t-1}$ a $O_t$),
tratando o ativo como se nada ocorresse fora do pregão. Para o sistema bancário brasileiro
no período 2019--2025, isso é particularmente problemático já que decisões do Comitê de Política
Monetária (Copom) são divulgadas após o fechamento do mercado, resultados trimestrais são
publicados \textit{after-market} e dados macroeconômicos norte-americanos chegam durante a
tarde-noite no horário de Brasília. A literatura empírica em ações brasileiras estima que
entre 40\% e 50\% da variância total ocorre \textit{overnight}, de modo que descartar essa
componente subestima sistematicamente a volatilidade efetiva. Adicionalmente, Garman-Klass
assume \textit{drift} intradiário nulo, o que o torna positivamente viesado em janelas com
tendência marcada como em março de 2020 (COVID-19) e janeiro de 2023 (Americanas). O
estimador de Parkinson utiliza apenas os extremos ($H_t$ e $L_t$), sendo estritamente
menos eficiente do que Garman-Klass, que incorpora $O_t$ e $C_t$.

A volatilidade realizada ($RV$) intradiária exigiria dados de \textit{tick},
indisponíveis abertamente para o período amostral completo. Mesmo com acesso a dados de
alta frequência, medições a 5 minutos produzem viés por \textit{bid-ask bounce}, e as
correções técnicas associadas fogem do escopo deste trabalho. A presença de \textit{mid
caps} na amostra agrava o problema de modo que dias sem negociação contínua geram zeros artificiais
que inflam $RV$. A abordagem GARCH(1,1) introduziria uma camada paramétrica adicional
antes do VAR, gerando duas etapas de modelagem encadeadas que complicam a interpretação
e a robustez, além de operar apenas com retornos \textit{close-to-close}, descartando
a informação intradiária que torna os estimadores \textit{range-based} mais eficientes.

\subsubsection*{Período amostral}

O período amostral abrange de fevereiro de 2019 a dezembro de 2025. Esse recorte garante
que todas as oito instituições já tenham concluído seus IPOs e apresentem histórico
suficiente de negociação. Adicionalmente, o período concentra os principais episódios de
estresse do sistema financeiro brasileiro recente, cujos efeitos sobre a volatilidade
bancária são centrais para a análise dinâmica: a crise da COVID-19 e os \textit{circuit breakers} de março de 2020, 
o ciclo de aperto monetário da Selic entre 2022 e 2023, a fraude contábil da Americanas (janeiro de
2023) e o estresse fiscal do segundo semestre de 2024. A Tabela~\ref{tab:eventos}
apresenta os principais eventos de \textit{drift} identificados na janela amostral.

\begin{table}[H]
\centering
\caption{Principais eventos de \textit{drift} no período amostral (2019-2025)}
\label{tab:eventos}
\small
\begin{tabular}{l>{\raggedright\arraybackslash}p{3.8cm}>{\raggedright\arraybackslash}p{6cm}}
\toprule
Período & Direção & Causa \\
\midrule
Mar./2020                & Negativo extremo               & \textit{Crash} COVID-19 \\
Abr.--dez./2020          & Positivo persistente           & Recuperação pós-\textit{crash} \\
Jan.--set./2022          & Negativo persistente           & Aperto da Selic \\
Jan.--mar./2023          & Negativo (ITUB4, BBDC4)        & Fraude Americanas \\
Ago.--dez./2023          & Positivo (\textit{large caps}) & Início do ciclo de corte da Selic \\
Mai./2024 (BRSR6)        & Negativo idiossincrático       & Enchentes no Rio Grande do Sul \\
Out.--dez./2024          & Negativo persistente           & Crise fiscal pós-pacote governamental \\
Jan.--abr./2025          & Negativo (sistema)             & Tarifas Trump pós-posse \\
Mai.--dez./2025          & Positivo (\textit{large caps}) & Expansão de \textit{spread} bancário \\
\bottomrule
\end{tabular}
\end{table}

\subsection*{6.1.3 Tratamento dos dados}

Os procedimentos de tratamento são aplicados na seguinte ordem.

\subsubsection*{1.\ Extração dos dados}

Os dados OHLC de todas as oito instituições foram obtidos da base ComDinheiro.
Os preços são fornecidos já ajustados retroativamente por proventos considerando dividendos, 
juros sobre capital próprio, bonificações, \textit{splits}, \textit{inplits} e cisões 
(\textit{spin-offs}). A escolha por essa fonte foi motivada por duas limitações identificadas.
Primeiro, o Banco Pan (BPAN4) deixou de ser negociado ao final de janeiro de 2026, 
em decorrência da incorporação ao BTG Pactual, o que tornou o \textit{ticker} indisponível 
em provedores que não mantêm o histórico de ativos descontinuados. Segundo, o ajuste automático
de provedores como o Yahoo Finance corrige dividendos e \textit{splits} mas não cisões. A série
de ITUB4, por exemplo, registraria uma queda de aproximadamente $-19,8\%$ decorrente
da cisão do braço de operações de cartões do Itaú, em um movimento que não reflete volatilidade
de mercado real. Caso não corrigido, a estimação dos \textit{spillovers} seria distorcida
nas janelas que cobrem o período.

\subsubsection*{2.\ Verificação de consistência}

Antes da estimação da volatilidade, os dados OHLC são submetidos a verificações de
consistência lógica: $H_t \geq L_t$, $H_t \geq O_t$, $H_t \geq C_t$, $L_t \leq O_t$,
$L_t \leq C_t$ e $O_t, C_t > 0$. Violações indicam erros de digitação ou de
\textit{parsing} no provedor de dados, que produzem valores \texttt{NaN} na estimação de
Yang-Zhang e falham silenciosamente.

\subsubsection*{3.\ Filtro de liquidez}

Dias com volume de negociação inferior a 1.000 ações ou volume financeiro inferior a
R\$\,100 mil são marcados como \texttt{NA}. Para as \textit{mid caps} ABCB4, BPAN4 e
BRSR6, dias com volume muito reduzido produzem \textit{ranges} ($H_t - L_t$) que refletem
o \textit{bid-ask spread} da microestrutura, e não a volatilidade efetiva do ativo.
Observações filtradas são tratadas na etapa de imputação.

\subsubsection*{4.\ Tratamento de \textit{outliers}}

Retornos extremos são identificados pelo escore padronizado
$z_t = (r_t - \bar{r})/\hat{\sigma}_r$ calculado sobre os retornos
\textit{close-to-close}, sendo $|z_t| > 5$ o critério de sinalização. Para cada
observação sinalizada, verifica-se: se múltiplas instituições registram retorno extremo na
mesma data, o evento é classificado como real e mantido; se apenas uma instituição
apresenta retorno extremo sem notícia identificável, a observação é conferida em fonte
alternativa. Como referência, 12 e 16 de março
de 2020 (\textit{circuit breakers} da COVID-19) e 12 de janeiro de 2023 (Americanas,
restrito a ITUB4 e BBDC4) são eventos reais e são preservados na amostra.

\subsubsection*{5.\ Estimação da volatilidade}

A volatilidade Yang-Zhang é estimada conforme as
equações~\eqref{eq:yz}--\eqref{eq:k}, com janela rolante de $n = 21$ dias úteis. Aplica-se
um piso numérico de $10^{-8}$ antes da transformação logarítmica:
$y_t = \log\bigl(\max(\hat{\sigma}_{\mathrm{YZ},t},\,10^{-8})\bigr)$. Os primeiros
20 dias de cada série recebem \texttt{NA} em função da janela \textit{rolling}.

\subsubsection*{6.\ Alinhamento e imputação}

As oito séries de log-volatilidade são alinhadas por interseção de datas de pregão,
formando um painel $T \times 8$. Dias em que todas as séries apresentam \texttt{NA}
(finais de semana, feriados) são excluídos. Observações isoladas com \texttt{NA}
decorrentes do filtro de liquidez ou de erros de provedor são imputadas por filtro de
Kalman aplicado a um modelo ARIMA selecionado automaticamente\footnote{O filtro de Kalman
fornece estimativas ótimas de estados não observados em um modelo de espaço de estados,
minimizando o erro quadrático médio da predição. Neste contexto, o modelo ARIMA descreve
a dinâmica autorregressiva da série de log-volatilidade e o filtro utiliza toda a
informação disponível, observações passadas e futuras (\textit{smoother}) para
imputar os valores ausentes. Para detalhes, consultar \textcite{moritz2017}.}.
A integridade do painel é verificada ao final pela ausência de qualquer \texttt{NA}
residual.

\subsection*{6.1.4 Propriedades estatísticas}

A estacionariedade de cada série de log-volatilidade é testada por meio do teste de
Dickey-Fuller Aumentado (ADF), o teste KPSS, o teste DF-GLS, 
o teste de Phillips-Perron e o teste de Zivot-Andrews. Séries de log-volatilidade financeira
frequentemente exibem memória longa, de modo que tanto o ADF quanto o KPSS podem rejeitar
suas respectivas hipóteses nulas simultaneamente \parencite{andersen2001, andersen2003, diebold2001}.
Esse padrão de rejeição simultânea não constitui evidência de raiz unitária. A estimação por
janela rolante também constitui uma resposta natural à não-estacionariedade de longo prazo. 
Ao estimar o VAR sobre janelas de $W = 200$ dias, a dinâmica de memória longa e as quebras
de regime da amostra completa são atenuadas e cada janela se aproxima de um processo estacionário
no curto prazo.
A presença de autocorrelação serial e heterocedasticidade
condicional é verificada pelos testes de Ljung-Box e ARCH-LM, respectivamente.
% Estatísticas descritivas detalhadas encontram-se no Apêndice Estatístico.

% ──────────────────────────────────────────────────────────────────────────────
\section{Topologia de Redes via Diebold-Yilmaz}
\label{subsec:dy}
% ──────────────────────────────────────────────────────────────────────────────

\subsection*{Especificação do VAR}

Seja $\mathbf{y}_t = (y_{1t}, \ldots, y_{Nt})'$ o vetor $N \times 1$ das séries de
log-volatilidade, com $N = 8$. O modelo VAR de ordem $p$ é:

\begin{equation}
  \mathbf{y}_t
    = \mathbf{c}
    + \sum_{i=1}^{p} \boldsymbol{\Phi}_i\,\mathbf{y}_{t-i}
    + \boldsymbol{\varepsilon}_t,
  \qquad
  \boldsymbol{\varepsilon}_t \sim (\mathbf{0},\,\boldsymbol{\Sigma}),
  \label{eq:var}
\end{equation}

\noindent onde $\mathbf{c}$ é um vetor de constantes, $\boldsymbol{\Phi}_i$ são matrizes
$N \times N$ de coeficientes autorregressivos e $\boldsymbol{\varepsilon}_t$ é um vetor
de inovações com matriz de covariância $\boldsymbol{\Sigma}$. Os parâmetros são estimados
por Mínimos Quadrados Ordinários (MQO) equação a equação, estimador implementado pelo pacote
\texttt{ConnectednessApproach} \parencite{gabauer2022}.

A especificação principal adota $p = 2$ defasagens. Essa escolha reflete dois critérios.
Primeiro, a parcimônia: com janela rolante de $W = 200$ dias úteis, o modelo VAR$(p)$
com $N = 8$ variáveis possui $N^2 p + N = 136$ parâmetros no total. Como a estimação
é feita por MQO equação a equação, o parâmetro relevante para a estabilidade numérica
não é o total do sistema, mas a razão observações/parâmetros por equação: cada uma das
$N$ equações possui $Np + 1 = 17$ regressores, de modo que essa razão é de
$T_w / (Np+1) = 200/17 \approx 11{,}7$ para $p=2$. Com $p = 4$, essa razão cai para
$200/33 \approx 6{,}1$, deteriorando a estabilidade numérica da estimação. 
Segundo, consistência com a literatura: \textcite{diebold2009} e \textcite{diebold2012} 
adotam $p = 2$ como especificação de referência em amostras similares. 
% A Tabela~\ref{tab:bic} apresenta o critério de informação de Schwarz (BIC) para 
% $p \in \{1, 2, 3, 4\}$, estimado sobre o \textit{full sample}. O Apêndice de Robustez 
% reporta a sensibilidade dos índices de spillover a diferentes valores de $p$.

% TODO (Mono 2): rodar VARselect() no full sample para gerar a Tabela~\ref{tab:bic}
% (BIC para p = 1,2,3,4) e descomentar o bloco abaixo.
%
% \begin{table}[H]
% \centering
% \caption{Critério de informação de Schwarz (BIC) por ordem de defasagem}
% \label{tab:bic}
% \begin{tabular}{cc}
% \toprule
% $p$ & BIC \\
% \midrule
% 1 & ... \\
% 2 & ... \\
% 3 & ... \\
% 4 & ... \\
% \bottomrule
% \end{tabular}
% \end{table}

A condição de estabilidade do VAR exige que todas as raízes do polinômio característico
$\det\!\bigl(\mathbf{I}_N - \boldsymbol{\Phi}_1 z - \cdots - \boldsymbol{\Phi}_p
z^p\bigr) = 0$ encontrem-se fora do círculo unitário no plano complexo,
equivalentemente, que os módulos dos autovalores da matriz companheira sejam estritamente
menores que um. Essa condição garante que o VAR admita a representação de média móvel de
ordem infinita $\mathbf{y}_t = \sum_{h=0}^{\infty} \mathbf{A}_h\,\boldsymbol{\varepsilon}_{t-h}$,
utilizada na decomposição de variância subsequente.

\subsection*{Identificação via GFEVD}

A identificação dos spillovers baseia-se na \textit{Generalized Forecast Error Variance
Decomposition} (GFEVD) de \textcite{pesaran1998}. Dado o horizonte de previsão $\tau$, a
fração da variância do erro de previsão de $i$ atribuída a choques em $j$ é:

\begin{equation}
  \theta^{\tau}_{ij}
    = \frac{
        \sigma_{jj}^{-1}
        \displaystyle\sum_{h=0}^{\tau-1}
        \bigl(\mathbf{e}_i'\,\mathbf{A}_h\,\boldsymbol{\Sigma}\,\mathbf{e}_j\bigr)^{2}
      }{
        \displaystyle\sum_{h=0}^{\tau-1}
        \mathbf{e}_i'\,\mathbf{A}_h\,\boldsymbol{\Sigma}\,\mathbf{A}_h'\,\mathbf{e}_i
      },
  \label{eq:gfevd}
\end{equation}

\noindent onde $\sigma_{jj}$ é o $j$-ésimo elemento diagonal de $\boldsymbol{\Sigma}$ e
$\mathbf{e}_i$ é o $i$-ésimo vetor canônico. A identificação generalizada de Pesaran e
Shin distingue-se da decomposição de Cholesky por não impor ordenação recursiva,
eliminando a sensibilidade dos resultados a permutações da ordem do VAR. Como
contrapartida, as linhas da matriz $\boldsymbol{\Theta}^{\tau} = [\theta^{\tau}_{ij}]$ não
somam necessariamente um. Aplica-se normalização linha-a-linha:

\begin{equation}
  \tilde{\theta}^{\tau}_{ij}
    = \frac{\theta^{\tau}_{ij}}{\displaystyle\sum_{j=1}^{N}\theta^{\tau}_{ij}},
  \label{eq:norm}
\end{equation}

\noindent de modo que $\sum_{j=1}^{N}\tilde{\theta}^{\tau}_{ij} = 1$ e
$\sum_{i,j=1}^{N}\tilde{\theta}^{\tau}_{ij} = N$ por construção. O horizonte é fixado em
$\tau = 10$ dias úteis, seguindo a convenção de
\textcite{diebold2009}, \textcite{diebold2012} e \textcite{diebold2014}.

\subsection*{Índices de \textit{spillover}}

A partir da matriz normalizada $\tilde{\boldsymbol{\Theta}}^{\tau}$, definem-se os índices
de conectividade de \textcite{diebold2012}: o índice de conectividade total, os índices
direcionais de transmissão e recepção de choques, o índice líquido por instituição e o
índice líquido par-a-par.

\begin{align}
  \text{TCI}
    &= \frac{1}{N}
       \sum_{\substack{i,j=1\\i\neq j}}^{N}
       \tilde{\theta}^{\tau}_{ij}
       \times 100,
  \label{eq:tci}\\[6pt]
  \text{FROM}_i
    &= \sum_{\substack{j=1\\j\neq i}}^{N}
       \tilde{\theta}^{\tau}_{ij}
       \times 100,
  \label{eq:from}\\[6pt]
  \text{TO}_j
    &= \sum_{\substack{i=1\\i\neq j}}^{N}
       \tilde{\theta}^{\tau}_{ij}
       \times 100,
  \label{eq:to}\\[6pt]
  \text{NET}_i
    &= \text{TO}_i - \text{FROM}_i,
  \label{eq:net}\\[6pt]
  \text{NET}_{ij}
    &= \tilde{\theta}^{\tau}_{ji} - \tilde{\theta}^{\tau}_{ij}.
  \label{eq:pairwise}
\end{align}

O \textit{Total Connectedness Index} (TCI) em~\eqref{eq:tci} mede a fração média da
variância do erro de previsão de cada banco explicada por choques nos demais.
FROM$_i$~\eqref{eq:from} quantifica a receptividade sistêmica de $i$; TO$_j$~\eqref{eq:to}
mede a transmissão sistêmica de $j$; NET$_i$~\eqref{eq:net} classifica as instituições
como transmissores líquidos (NET$_i > 0$) ou receptores líquidos (NET$_i < 0$);
NET$_{ij}$~\eqref{eq:pairwise} captura a direção líquida do fluxo de risco entre cada par
de instituições.

Para além dos índices direcionais propostos por \textcite{diebold2012}, este trabalho
incorpora um conjunto de métricas de teoria de redes, 
centralidade de intermediação (\textit{betweenness}), centralidade de autovetor, assortatividade e
PageRank, calculadas sobre o grafo de conectividade $G = (V, E, W)$ definido na
Seção~\ref{subsec:redes}. Essa incorporação constitui uma contribuição adicional do
trabalho. Enquanto a literatura de \textit{spillover} de volatilidade concentra-se nos
índices direcionais de primeira ordem (quanto cada instituição transmite e recebe
diretamente), as métricas de rede capturam a posição estrutural de cada banco na
topologia de contágio, incluindo efeitos indiretos e recursivos. O caso mais ilustrativo
é o PageRank, cuja formulação recursiva atribui a $i$ um valor elevado quando $i$ recebe 
\textit{spillovers} de instituições que, por sua vez, também possuem PageRank elevado.
A interpretação econômica é a de contribuição recursiva de risco, de modo que um banco 
é sistemicamente importante não apenas pela intensidade de seus \textit{spillovers} diretos, 
mas por estar conectado a outras instituições igualmente centrais, lógica análoga à do 
\textit{framework} de bancos sistemicamente importantes (D-SIB) adotado pelo Banco Central 
do Brasil. O contraste central deste trabalho é a confrontação entre o ranking de 
importância sistêmica gerado por essas métricas de rede e a segmentação oficial de D-SIBs 
publicada pelo BCB.

\subsection*{Análise dinâmica}

Os índices de \textit{spillover} são estimados em janela rolante de $W = 200$ dias úteis
($\approx 10$ meses de negociação). Esse valor representa um equilíbrio entre a
estabilidade numérica da estimação e a responsividade a mudanças de regime. O horizonte
$\tau = 10$ segue a convenção da literatura de referência
\parencite{diebold2009,diebold2012,diebold2014,demirer2018,antonakakis2020}.

%O Apêndice de Robustez reporta a sensibilidade dos resultados a $p \in \{1,2,3,4\}$,
%$W \in \{150,200,250\}$ dias úteis e $\tau \in \{5,10,20\}$ dias, além da comparação entre
%spillovers de log-volatilidade e de retorno e da estabilidade dos rankings em
%sub-amostras pré- e pós-pandemia.

% ──────────────────────────────────────────────────────────────────────────────
\section{Visualização de Redes}
\label{subsec:redes}
% ──────────────────────────────────────────────────────────────────────────────

A matriz normalizada $\tilde{\boldsymbol{\Theta}}^{\tau}$ define a matriz de adjacência de um
grafo dirigido ponderado $G = (V, E, W)$, em que $V$ representa as oito instituições, $E$
o conjunto de arestas de transmissão de variância e $W$ os pesos $\tilde{\theta}^{\tau}_{ij}$.
A convenção de direção é: $\tilde{\theta}^{\tau}_{ij}$ representa o fluxo de $j$ para $i$
(o choque em $j$ explica parcela da variância do erro de previsão de $i$). Como o pacote
\texttt{igraph} interpreta o elemento $\mathrm{mat}[i,j]$ como aresta saindo de $i$ e
chegando em $j$, a matriz é transposta antes da construção do grafo para preservar a
semântica econômica.

Os índices definidos na Subseção~\ref{subsec:dy} são mapeados em atributos visuais do
grafo: (i)~o tamanho do nó é proporcional ao índice TO$_i$; (ii)~a cor do nó é
codificada pelo índice NET$_i$ por meio de uma paleta divergente, distinguindo
transmissores líquidos de receptores líquidos; (iii)~a espessura de cada aresta é
proporcional ao \textit{spillover} par-a-par NET$_{ij}$; e (iv)~a forma ou contorno do nó
codifica o tipo de controle institucional.

A rede DY é quase completa: com $N = 8$, há até 56 arestas dirigidas. Para fins de
visualização exclusivamente, aplica-se um \textit{threshold} de 5\% sobre os pesos
$\tilde{\theta}^{\tau}_{ij}$, retendo apenas as conexões mais intensas. Esse filtro não afeta
a inferência conduzida a partir das tabelas e séries temporais dos índices DY.

O grafo do \textit{full sample} é posicionado pelo algoritmo de Fruchterman-Reingold. Para
os painéis comparativos entre subperíodos (pré-COVID, COVID-19, aperto da Selic e crise
Americanas), as coordenadas dos nós são calculadas uma única vez sobre o \textit{full
sample} e mantidas fixas em todos os painéis. Essa escolha é metodológica, garantindo que as
diferenças visuais entre painéis reflitam reconfiguração da topologia da rede, e não
reposicionamento dos nós.

A representação em grafo é um instrumento heurístico e comunicativo. As escolhas de
\textit{threshold}, algoritmo de \textit{layout} e paleta de cores afetam a percepção
visual, mas não os coeficientes estimados. Toda a inferência sobre transmissão de risco
entre as instituições é conduzida a partir das tabelas e das séries temporais dos índices
DY e das métricas de centralidade definidos na Subseção~\ref{subsec:dy}.

\chapter{Resultados}
\label{cap:resultados}

\postextual

% ---- Referências bibliográficas ----

% No final do documento, antes de \end{document}:
% \bibliography{referencias}

\printbibliography

\end{document}